% Lösungen Kenngrößen von Verteilungen (zusammenbau v0.4, Heft)
\einheitenkopf{Kenngrößen von Verteilungen}

\erg{27}{a) Auszahlung \quad b) fehlend $1 - 0{,}5 - 0{,}3 = 0{,}2$; $E(X) = 0 \cdot 0{,}5 + 1 \cdot 0{,}3 + 2 \cdot 0{,}2 = 0{,}7$ \quad c) $E(X) = 50 \cdot 0{,}02 + 5 \cdot 0{,}1 + 0 = 1{,}50$ (Euro)}
\erg{28}{a) $E(\text{Auszahlung}) = 0{,}25 \cdot 10 = 2{,}50$ (Euro), weniger als der Einsatz; Gewinn des Anbieters $= 3 - 2{,}50 = 0{,}50$ (Euro) je Spiel \quad b) Nein: $E(\text{Auszahlung}) = 0 \cdot 0{,}5 + 5 \cdot 0{,}3 + 12 \cdot 0{,}2 = 3{,}90$ (Euro) liegt unter dem Einsatz von 4\,€; Jana verliert im Mittel 0{,}10\,€ je Spiel, der Anbieter gewinnt so viel \quad c) Der zweite Summand gehört zu genau einer roten Kugel: 4 ist die Auszahlung in Euro, der Faktor 2 zählt die zwei Reihenfolgen rot–blau und blau–rot, $0{,}3 \cdot 0{,}7$ ist die Wahrscheinlichkeit eines solchen Pfades \quad d) Der zweite Summand gehört zu genau einmal Grün: 1 ist die Auszahlung in Euro, die 2 zählt die zwei Pfade grün–nicht grün und nicht grün–grün, $\frac{1}{4} \cdot \frac{3}{4}$ ist die Wahrscheinlichkeit eines Pfades; die 2 ist eine Anzahl von Pfaden, kein Betrag}
\erg{29}{a) Erwartungswert $300 \cdot \frac{1}{6} = 50$; Abweichung $\frac{50 - 42}{50} = 16\,\%$ unter dem Erwartungswert}
\erg{30}{a) Erwartungswert $150 \cdot 0{,}06 = 9$; Abweichung $\frac{12 - 9}{9} \approx 33{,}3\,\%$ über dem Erwartungswert}
\erg{31}{a) vorwärts \quad b) $0{,}25 \cdot a = 2$, also $a = 8$ (Euro) \quad c) $\frac{4}{6} \cdot 1 + \frac{2}{6} \cdot b = 3$, also $4 + 2b = 18$ und $b = 7$}
\erg{32}{a) $\frac{6}{6 + x} \cdot \frac{5}{5 + x} = \frac{1}{3}$ führt auf $(6 + x)(5 + x) = 90$, also $x^2 + 11x - 60 = 0$; $x_1 = 4$, $x_2 = -15$ entfällt; also $x = 4$ \quad b) $3 \cdot \left(3p + \frac{1 - p}{3}\right)^2 = 3$, also $3p + \frac{1 - p}{3} = 1$ oder $= -1$; $p = \frac{1}{4}$ ($p = -\frac{1}{2}$ entfällt); Gold : Silber $= 1 : 3$}
\erg{33}{a) $4q^2 + 2 \cdot 2q(1 - q) + (1 - q)^2 = 2{,}25$ ergibt $q^2 + 2q + 1 = 2{,}25$, mal 4 die Gleichung; $q = \frac{1}{2}$ ($q = -\frac{5}{2}$ entfällt)}
\erg{34}{a) $16q^2 + 4 \cdot 2q(1 - q) + (1 - q)^2 = 4$ ergibt $9q^2 + 6q - 3 = 0$, geteilt durch 3 die Gleichung; $q = \frac{1}{3}$ ($q = -1$ entfällt)}
\erg{35}{a) mit $P(X = 2) = 0{,}8 - a$: $1 \cdot a + 2 \cdot (0{,}8 - a) + 4 \cdot 0{,}2 = 2{,}2$, also $P(X = 1) = 0{,}2$ und $P(X = 2) = 0{,}6$}
\erg{36}{a) je Tag im Mittel $270 : 30 = 9$ Münzen; mit $P(10) = 0{,}95 - p$: $5p + 10 \cdot (0{,}95 - p) + 40 \cdot 0{,}05 = 9$, also $p = P(5) = 0{,}5$ und $P(10) = 0{,}45$}
\erg{37}{a) Binomialformel – Bernoulli-Kette mit fester Trefferwahrscheinlichkeit \quad b) $p = \frac{50}{100} = 0{,}5$; $\sigma = \sqrt{100 \cdot 0{,}5 \cdot 0{,}5} = 5$ \quad c) symmetrisch heißt $p = 0{,}5$; $n = 7 : 0{,}5 = 14$}
\erg{38}{a) $12 \cdot (1 - p) = 3$, also $1 - p = 0{,}25$ und $p = 0{,}75$; $n = 12 : 0{,}75 = 16$ \quad b) $\sigma^2 = \mu \cdot (1 - p)$: $9 = 12 \cdot (1 - p)$, also $p = 0{,}25$; $n = 12 : 0{,}25 = 48$; $P(X = 13) = \frac{48!}{13! \cdot 35!} \cdot 0{,}25^{13} \cdot 0{,}75^{35}$ \quad c) $18 \cdot p \cdot (1 - p) = 6{,}25$ verlangt $p \cdot (1 - p) = 6{,}25 : 18 \approx 0{,}35$; das Produkt ist als nach unten geöffnete Parabel mit Scheitel bei $p = 0{,}5$ höchstens 0{,}25, also gibt es kein solches $p$ \quad d) Der Bogen reicht von $p = 0$ bis $p = 1$, also $1 : 20 = 0{,}05$ je Kästchen; bei $p = 0{,}25$ ist $\sigma = \sqrt{1\,200 \cdot 0{,}25 \cdot 0{,}75} = 15$, fünf Kästchen, also 3 je Kästchen; das Maximum $\sqrt{300} \approx 17{,}3$ liegt auf keiner Gitterlinie und taugt nicht zum Skalieren \quad e) $p$-Achse: $1 : 10 = 0{,}1$ je Kästchen; bei $p = 0{,}1$ ist $\sigma = \sqrt{3\,600 \cdot 0{,}1 \cdot 0{,}9} = 18$, vier Kästchen, also $18 : 4 = 4{,}5$ je Kästchen; der höchste Punkt $\sigma = 30$ liegt bei $6{,}\overline{6}$ Kästchen, zwischen zwei Gitterlinien}
\erg{39}{a) $\sqrt{n \cdot 0{,}5 \cdot 0{,}5} = 5$, quadriert $0{,}25 \cdot n = 25$, also $n = 100$}
\erg{40}{a) $\sqrt{n \cdot 0{,}1 \cdot 0{,}9} = 6$, quadriert $0{,}09 \cdot n = 36$, also $n = 400$}
\erg{41}{a) Er ist der Durchschnitt der Trefferzahlen vieler Durchgänge. \quad b) höchste Säule bei 3, also $\mu = 3$; $p = 3 : 10 = 0{,}3$ \quad c) $X$: höchste Säule bei 2, $p_X = 0{,}2$; $Y$: höchste Säule bei 6, $p_Y = 0{,}6$; $Y$ hat die größere Trefferwahrscheinlichkeit, obwohl ihre höchste Säule niedriger ist}
\erg{42}{a) höchste Säule bei 6, also $\mu = 30 \cdot p = 6$ und $p = 0{,}2$ (nicht die Säulenhöhe ablesen) \quad b) Das Maximum liegt zwischen 7 und 8, der Erwartungswert ist also $\mu = 7{,}5 = n \cdot 0{,}5$ und nicht ganzzahlig; bei geradem $n$ wäre $n \cdot 0{,}5$ ganzzahlig, also ist $n = 15$ ungerade}
